%==============================================================================
\section{Related Work}
\label{sec:related}

\subsection{Convolutional segmentation networks}
\label{sec:related-cnn}

U-Net \citep{ronneberger2015} established the encoder--decoder
architecture with skip connections that remains the dominant template for
medical image segmentation. Subsequent work has refined the encoder
(residual blocks / ResNet \citep{he2016} and later compound-scaled
backbones), the decoder
(transposed convolution vs.\ bilinear upsampling, dense connections
\citep{huang2017}), and the skip mechanism
(attention gates \citep{oktay2018}, dense skips). Squeeze-and-excitation
blocks \citep{hu2018} provide channel-wise recalibration. Atrous spatial
pyramid pooling (ASPP) in DeepLabV3+ \citep{chen2018encoder} replaces the
deep encoder with dilated convolutions to enlarge the receptive field
without down-sampling. Feature Pyramid Networks
\citep{lin2017} aggregate features across multiple scales for dense
prediction.

The nnU-Net framework \citep{isensee2021} integrates these architectural
choices with automatic data-dependent hyperparameter selection (target
spacing, patch size, batch size, augmentation pipeline) and has set
strong baselines on numerous medical benchmarks including CAMUS.

\subsection{Transformer-based segmentation networks}
\label{sec:related-trans}

Self-attention's quadratic complexity made global attention impractical
for high-resolution images until Vision Transformers
\citep{dosovitskiy2021} demonstrated patch-based tokenisation. TransUNet
\citep{chen2021transunet} was the first work to bring this design to
medical segmentation, combining a CNN encoder (typically ResNetV2-BiT
\citep{kolesnikov2020}) with a ViT-B/16 bottleneck and a cascaded
upsampler decoder. The hybrid design retains the inductive bias of
convolutions at high resolution while granting global context at the
bottleneck.

Swin Transformer \citep{liu2021swin} introduced windowed self-attention
with shifted windows, restoring linear complexity in image size. Swin-UNet
\citep{cao2022} adapts this to segmentation with a fully Transformer-based
U-Net. SwinV2 \citep{liu2022swinv2} and Swin-MAE
\citep{baochen2023swinmae} further refine pretraining strategies.

\subsection{Echocardiographic segmentation and CAMUS}
\label{sec:related-camus}

The CAMUS dataset \citep{leclerc2019} (\emph{Cardiac Acquisitions for
Multi-structure Ultrasound Segmentation}) provides 2D apical-2-chamber and
apical-4-chamber acquisitions from 500 patients, each annotated by an
expert echocardiographer at end-diastole and end-systole with four
classes (background, LV-endo, LV-epi, LA). The official split assigns
$400$ patients to training, $50$ to validation, and $50$ to test, with
test annotations held out by the challenge organisers. Image quality is
graded into \emph{Good}, \emph{Medium}, and \emph{Poor} categories by
the same experts.

The original CAMUS paper reported nnU-Net achieving
HD $\approx 4.3$~mm on LV-endo and $\sim$$9\%$ EF MAE.
Subsequent work \citep{painchaud2022, smistad2021}
has reported incremental improvements using deeper architectures, custom
losses, and view-aware training. Several papers have additionally
explored sequence-aware models on CAMUS half-cycle annotations
\citep{painchaud2022}, but these are out of scope for the present
benchmark.

Activity on echocardiographic segmentation has continued along four
largely separate lines. The first pursues \emph{architectural} refinement:
bi-directional token matching between frames \citep{liu2025botm},
dynamic-learning U-Nets that enforce temporal consistency across the
cardiac cycle \citep{qu2025dylunet}, contour-guided query-based fusion
aimed explicitly at boundary quality \citep{ullah2026contour}, and
selective state space models applied to paediatric echocardiography
\citep{ye2024pmamba}. The second addresses \emph{domain shift}, which is
the dominant obstacle to deploying any of these models across scanners
and sites: reinforcement-learning-based domain adaptation
\citep{judge2024domain, judge2025rlda} and interpretable quantification of
the shift itself \citep{elyasi2026domainshift}. The third applies
\emph{foundation models} and self-supervision, either by adapting
general-purpose promptable segmenters to medical images
\citep{ma2024samed}, by generating synthetic cardiac ultrasound for
augmentation \citep{tiago2024synthetic, reynaud2025echoflow}, or by asking
what self-supervised echocardiographic representations actually transfer
\citep{majchrowska2026ssl}. The fourth targets the \emph{clinical
endpoint} directly, predicting or explaining ejection fraction rather than
the mask that produces it \citep{zhang2024lvpredict, thomas2024echonarrator,
li2025echovlm, wang2026echoagent}.

What these lines share is a comparison protocol inherited from whichever
baseline was convenient. Each reports its own training recipe, its own
resolution, and — as we show in Section~\ref{sec:results-boundary} — not
always the same definition of a millimetre. That is the gap this benchmark
addresses.

EchoNet-Dynamic \citep{ouyang2020} provides a complementary large-scale
video dataset for EF regression. We focus on CAMUS as the standard
benchmark for \emph{segmentation-derived} EF. Because EchoNet-Dynamic
annotates only the left-ventricular endocardium in a single apical view,
it cannot exercise the LV-epicardium, left-atrium, or biplane-EF
components of our protocol; whether the architectural trade-offs reported
here transfer to other echocardiographic cohorts is therefore a question
we leave to future cross-dataset work rather than a claim we make here.

\subsection{Validation practice and metric reporting}
\label{sec:related-validation}

Running alongside the architectural literature is a smaller body of work
arguing that the field's \emph{measurement} practice, not its model
design, is now the limiting factor. \citet{reinke2021limitations} catalogue
the ways standard segmentation metrics mislead when applied without regard
to structure size, shape, or resolution, and the Metrics Reloaded
framework \citep{maierhein2024metrics} turns that catalogue into
problem-driven recommendations for choosing and reporting them.
\citet{isensee2024revisited} show that much of the reported progress in 3D
medical segmentation dissolves under rigorous validation, with a
well-configured U-Net still matching or beating architectures claimed to
supersede it. \citet{kazaj2025claims} reach a similar conclusion for the
CNN-versus-Transformer-versus-Mamba question specifically, and
\citet{urooj2026mirage} audit a segmentation subfield and find that
reported rankings shift once evaluation inconsistencies are removed.

Boundary metrics are especially exposed. Hausdorff distances are reported
in millimetres, but a millimetre is only defined once the grid on which
the distance transform is computed and the physical spacing attached to
that grid refer to the same image. Networks are trained at a fixed
resolution while clinical images are not, so the two can silently diverge.
We quantify the consequence on CAMUS in
Section~\ref{sec:results-boundary}: the same predictions score
$1.8$~mm or $4.1$~mm depending on which convention is used, and only the
latter is comparable to the published baseline.

\subsection{What is missing}
\label{sec:related-gap}

A small number of recent works compare a handful of architectures on
CAMUS in their experimental sections, but none provide:
\begin{itemize}[leftmargin=*,nosep]
\item A unified training protocol across CNN, hybrid, and Transformer
      paradigms (papers typically inherit each architecture's own
      published hyperparameters, confounding architecture with
      protocol).
\item Per-class HD95 in millimetres with per-sample NIfTI spacing
      (most papers report Dice only, or use pixel-space HD95).
\item ED- versus ES-stratified reporting (essential for biplane
      Simpson's EF analysis).
\item Robustness analysis by CAMUS image-quality grade.
\item Pairwise Wilcoxon significance with Bonferroni correction and a
      critical-difference diagram.
\item Released model checkpoints and per-patient predictions for
      independent reanalysis.
\end{itemize}
This paper aims to fill each of these gaps.
